\section{Kubernetes, RKE2, and Containers}\label{sec:k8s}

\subsection{RKE2 Control Plane}

RKE2 runs a three-node HA control plane on VMs 201/202/203 (pve3/pve4/pve5,
\code{.51-.53}), all live-confirmed running (E-06). The API is fronted by a kube-vip VIP
at \code{192.168.10.54:6443}; the CNI is Cilium and load balancing is MetalLB over the
\code{.71-.75} pool. etcd runs as an embedded quorum across the three control-plane
nodes. The cluster is intentionally CPU-only: GPU scheduling is deferred because SLURM
and Ollama own the cards, and running RKE2 GPU scheduling on the same single card would
contend for exclusive device access.

Randy joined as a bare-metal storage worker with kubelet reservations sized to measured
hardware and a \code{node.netframe.io/\allowbreak role=storage:\allowbreak NoSchedule} taint, so only tolerating
infrastructure DaemonSets land there and ZFS/PBS are not cgroup-throttled. Persistent
storage is an \code{nfs-client} StorageClass backed by Randy \code{datastore/k8s}. A
documented operational detail worth preserving: agents register on the supervisor port
\code{:9345}, not the apiserver \code{:6443}.

\subsection{Container Gaps}

\begin{itemize}
\item \sevmed No NetworkPolicies: the pod network is flat with no default-deny
(QL-R15 / K8S-02).
\item \sevmed No Pod Security Admission (K8S-03), and Kubernetes secrets are likely
unencrypted at rest (K8S-04).
\item \sevmed The private registry is anonymous and a single replica (QL-R14 / K8S-01);
its TLS is provided by step-ca, and a \code{RegistryCertExpiringSoon} alert exists.
\item \sevlow Image pinning and floating \code{latest} tags were not exhaustively
audited; pinning discipline should be confirmed as workloads grow.
\item \sevinfo etcd is protected only at the VM-backup level; a dedicated etcd snapshot
schedule would speed control-plane recovery.
\end{itemize}

The recommended treatment (QL-REC15) is to add default-deny NetworkPolicies, Pod
Security Admission, secrets encryption, and registry authentication before the cluster
carries meaningful workloads. None of this should be applied without a maintenance window
and rollback.

\subsection{LXC and Docker Estate}

The operator convention is one service per LXC and "no Docker unless explicitly
requested," which keeps blast radius small and is a genuine strength. The running guests
were enumerated live (E-06): NPM (CT101), Vaultwarden (CT102), Grafana stack (CT103 on
pve4), Headscale (CT105), Homepage (CT106), Open WebUI (CT107), secondary Pi-hole
(CT108), plus VMs for OPNsense, Wazuh, Home Assistant, and the three RKE2 control-plane
nodes. Docker appears only inside a few LXCs where it is warranted: the Grafana stack
(Grafana, Prometheus, Loki, InfluxDB, Scrutiny) in CT103, Vaultwarden in CT102, and the
Homepage dashboard in CT106.

All guests are set \code{onboot=1}. Compose-versus-deployed drift could not be fully
verified read-only, since that would require exec into each container; it should be
spot-checked during normal maintenance. Consistent with the read-only mandate, this
review did not pull images, recreate containers, or restart any workload.

\begin{callout}
The one duplicate to resolve is CT104 on pve1, an old homepage instance that
documentation says was retired on 2026-06-24 but that is still running (DR-11 / R-12).
Confirm it is the retired skeleton, then stop and disable it (QL-REC17).
\end{callout}
